\subsection{Homotopic continuous maps}

DEFINITION 1. --- Let X and Y be topological spaces and let f and g be continuous maps from X to Y. A homotopy connecting f to g is a continuous map \(\sigma: X \times I \to Y\) such that, for every \(x \in X\), we have \(\sigma(x,0) = f(x)\) and \(\sigma(x,1) = g(x)\). We say that f is homotopic to g if there exists a homotopy connecting f to g.

We say that \(f\) is the origin and \(g\) the term of the homotopy \(\sigma\) (cf. below, III, p. 257, remark 2).

Let A be a subset of X. We say that the homotopy \(\sigma\) is fixed on A if, for every \(a \in \mathrm{A}\), the map \(t \mapsto \sigma(a, t)\) from I to Y is constant. In this case, the origin and the term of \(\sigma\) coincide at every point of A.

Let X and Y be topological spaces. We say that homotopies \(\sigma\colon X\times I\to Y\) and \(\tau\colon X\times I\to Y\) are juxtaposable if the term of \(\sigma\) is the origin of \(\tau\), in other words if one has \(\sigma(x,1)=\tau(x,0)\) for all \(x\in X\). In this case, the map \(\sigma*\tau\) from \(X\times I\) to Y defined by

\[
(\sigma * \tau) (x, t) = \left\{ \begin{array}{l l} \sigma (x, 2 t) & \text {for} 0 \leqslant t \leqslant 1 / 2 \\ \tau (x, 2 t - 1) & \text {for} 1 / 2 \leqslant t \leqslant 1 \end{array} \right.\tag{1}
\]

is continuous (TG, I, p. 19, prop. 4) and is a homotopy connecting the origin of \(\sigma\) to the term of \(\tau\). It is called the juxtaposed homotopy of the homotopies \(\sigma\) and \(\tau\).

If \(\sigma\colon X\times\mathbf{I}\to Y\) is a homotopy, the map \(\overline{\sigma}\colon X\times\mathbf{I}\to Y\) defined by \((x,t)\mapsto\sigma(x,1-t)\) is a homotopy connecting the term of \(\sigma\) to the origin of \(\sigma\). We have \(\overline{\overline{\sigma}}=\sigma\). If \(\sigma\) and \(\tau\) are juxtaposable homotopies from \(X\times\mathbf{I}\) to Y, the homotopies \(\overline{\tau}\) and \(\overline{\sigma}\) are juxtaposable, and one has \(\overline{\sigma*\overline{\tau}}=\overline{\tau}*\overline{\sigma}\).

PROPOSITION 1. --- Let X and Y be topological spaces. The relation "f is homotopic to g" is an equivalence relation on the set \(\mathcal{C}(\mathrm{X};\mathrm{Y})\) of continuous maps from X to Y.

Let \(f\) be an element of \(\mathcal{C}(\mathrm{X};\mathrm{Y})\). The map \(f\circ\mathrm{pr}_{1}\colon\mathrm{X}\times\mathrm{I}\to\mathrm{Y}\) is a homotopy connecting \(f\) to \(f\); this relation is therefore reflexive.

Let \(f\) and \(g\) be elements of \(\mathcal{C}(\mathrm{X};\mathrm{Y})\) and \(\sigma\colon\mathrm{X}\times\mathbf{I}\to\mathrm{Y}\) a homotopy connecting \(f\) to \(g\). The map \(\overline{\sigma}\colon\mathrm{X}\times\mathbf{I}\to\mathrm{Y}\) is then a homotopy connecting \(g\) to \(f\); the relation under consideration is therefore symmetric.

Let us finally prove that it is transitive. If \(f\), \(g\), and \(h\) are elements of \(\mathcal{C}(\mathrm{X};\mathrm{Y})\), \(\sigma\) a homotopy connecting \(f\) to \(g\), and \(\tau\) a homotopy connecting \(g\) to \(h\), then \(\sigma\) and \(\tau\) are juxtaposable and \(\sigma*\tau\) is a homotopy connecting \(f\) to \(h\).

Let X and Y be topological spaces. The equivalence relation “f is homotopic to g” in \(\mathcal{C}(\mathrm{X};\mathrm{Y})\) (prop. 1) is called the homotopy relation. The quotient set of \(\mathcal{C}(\mathrm{X};\mathrm{Y})\) by this relation is denoted {[}X;Y{]}. Its elements are called homotopy classes of continuous maps from X to Y. The homotopy class of a continuous map \(f\colon \mathrm{X}\to \mathrm{Y}\) will often be denoted {[}f{]}.

PROPOSITION 2. — Let X, Y, and Z be topological spaces, f and \(f'\) be continuous maps from X to Y, and g and \(g'\) continuous maps from Y to Z. If f is homotopic to \(f'\) and if g is homotopic to \(g'\), then \(g \circ f\) is homotopic to \(g' \circ f'\).

Let \(\sigma\) be a homotopy connecting \(f\) to \(f'\) and \(\tau\) a homotopy connecting \(g\) to \(g'\). Then, the map \(\theta\colon X\times I\to Z\) defined by \(\theta(x,t)=\tau(\sigma(x,t),t)\) is a homotopy connecting \(g\circ f\) to \(g'\circ f'\).

Let X, Y, and Z be topological spaces. Given homotopy classes \(\varphi \in [\mathrm{X};\mathrm{Y}]\) and \(\psi \in [\mathrm{Y};\mathrm{Z}]\), the maps \(g \circ f: \mathrm{X} \to \mathrm{Z}\), where \(f \in \varphi\), \(g \in \psi\), all belong to the same homotopy class (prop. 2), which is denoted \(\psi \circ \varphi\) and is called the composite homotopy class of the classes \(\psi\) and \(\varphi\). The map {[}X;Y{]} \(\times\) {[}Y;Z{]} \(\to\) {[}X;Z{]} which associates \(\psi \circ \varphi\) to \((\varphi,\psi)\) is called the composition map.

Let \(\varphi\in[\mathrm{X};\mathrm{Y}]\); we have \(\varphi=\varphi\circ[\mathrm{Id}_{\mathrm{X}}]=[\mathrm{Id}_{\mathrm{Y}}]\circ\varphi\).

Let X, Y, Z, and T be topological spaces, and let \(\varphi \in [\mathrm{X};\mathrm{Y}]\), \(\psi \in [\mathrm{Y};\mathrm{Z}]\), and \(\chi \in [\mathrm{Z};\mathrm{T}]\). The homotopy class \(\chi \circ (\psi \circ \varphi)\) is equal to \((\chi \circ \psi) \circ \varphi\); it is denoted \(\chi \circ \psi \circ \varphi\).

PROPOSITION 3. --- Let X be a topological space and let \((\mathrm{Y}_{j})_{j\in\mathrm{J}}\) be a family of topological spaces. The map from \([X;\prod_{j\in\mathrm{J}}Y_{j}]\) to the product set \(\prod_{j\in\mathrm{J}}[X;Y_{j}]\) defined by \(\varphi\mapsto([\mathrm{pr}_{j}]\circ\varphi)_{j\in\mathrm{J}}\) is bijective.

Surjectivity follows immediately from I, p. 25, prop. 1. Let us prove injectivity. Let \(f\) and \(g\) be continuous maps from X to \(\prod_{j \in \mathrm{J}} \mathrm{Y}_j\). For all \(j \in \mathrm{J}\), set \(f_j = \mathrm{pr}_j \circ f\) and \(g_j = \mathrm{pr}_j \circ g\); suppose \(f_j\) is homotopic to \(g_j\) and let \(\sigma_j\) be a homotopy connecting \(f_j\) to \(g_j\). The map \(\sigma = (\sigma_j)\) from \(\mathrm{X} \times \mathbf{I}\) to \(\prod_{j \in \mathrm{J}} \mathrm{Y}_j\) is continuous (loc. cit.); this is a homotopy connecting \(f\) to \(g\), whence the proposition.

COROLLARY. --- Let \((\mathrm{X}_{j})_{j\in\mathrm{J}}\) and \((\mathrm{Y}_{j})_{j\in\mathrm{J}}\) be families of topological spaces with the same index set. For every \(j\in J\), let \(f_{j}\) and \(g_{j}\) be continuous maps from \(X_{j}\) to \(Y_{j}\). If, for every \(j\in J\), the maps \(f_{j}\) and \(g_{j}\) are homotopic, then so are the product maps \(f\colon(x_{j})\mapsto(f_{j}(x_{j}))\) and \(g\colon(x_{j})\mapsto(g_{j}(x_{j}))\) from \(\prod_{j\in J}X_{j}\) to \(\prod_{j\in J}Y_{j}\).

Suppose that we have \([f_j] = [g_j]\) for all \(j \in \mathbf{J}\). We have \([\mathrm{pr}_j] \circ [f] = [f_j \circ \mathrm{pr}_j]\) and \([\mathrm{pr}_j] \circ [g] = [g_j \circ \mathrm{pr}_j]\); therefore \([\mathrm{pr}_j] \circ [f] = [\mathrm{pr}_j] \circ [g]\) by prop. 2, whence \([f] = [g]\) by prop. 3.

